\documentclass[10pt]{article}

\usepackage[utf8]{inputenc}

\usepackage[T1]{fontenc}

\usepackage{amsmath}

\usepackage{amsfonts}

\usepackage{amssymb}

\usepackage[version=4]{mhchem}

\usepackage{stmaryrd}

\usepackage{bbold}

\usepackage{graphicx}

\usepackage[export]{adjustbox}

\graphicspath{ {./images/} }



\begin{document}

спектра $\sigma(T)$, т. е. финитная $C^{1}$-функция в $\mathbb{C}$, для которой



$f=\left.\tilde{f}\right|_{\sigma(T)},\left|\frac{\partial \tilde{f}}{\partial \bar{z}}(\lambda)\right| \leqslant \operatorname{const} \cdot h_{\left\{M_{k}\right\}}(c \operatorname{dist}(\lambda, K)) ;$



вдесь



\[

\begin{gathered}

\frac{\partial \tilde{f}}{\partial \bar{z}}=\frac{1}{2}\left(\frac{\partial \tilde{f}}{\partial x}+i \frac{\partial \tilde{f}}{\partial y}\right), \\

h_{\left\{M_{k}\right\}}(r) \stackrel{\text { def }}{=} \inf _{n} r^{n-1} \frac{M_{n}}{n !},

\end{gathered}

\]



а оператор $T$ удовлетворяет условию



\[

\|R(\lambda, T)\| \leqslant\left({ }^{\{} M_{k}\right\}\left(\operatorname{dist}(\lambda, K)|\ln \operatorname{dist}(\lambda, K)|^{-1}\right) .

\]



C другой стороны, более широкие (чем $\operatorname{Hol}(\sigma(T))$ исчисления возникают как следствия оценок операторных многочленов $p(T)$; напр., если $X$ - гильбертово пространство, то неравенство Неймана - Хайнца



\[

\|p(T)\| \leqslant \max \{|p(\zeta)|:|\zeta| \leqslant\|T\|\}

\]



приводит к Ф. и. Сёкефальви-Надя-Фояша $(A$-алгебпа всех голоморфных и ограниченных в круге $\{\zeta \in$ $\in \mathbb{C}:|\zeta|<1\}$ функций, $T$ - сжатие без унитарных частей), а оно - $\mathrm{\kappa}$ многочисленным приложениям в теории функциональной модели для сжимающих операторов. Аналог неравенств Неймана - Хайнца для симметричных пространств функций порождает Ф. и. в терминах мультипликаторов (соответствующих сверточных пространств, [8]).



П р и м е н е ни я. Тип Ф. и., допускаемый оператором $T$, является инвариантом относительно линейного подобия $T \rightarrow V^{-1} T V$ и успешно используется для классификации операторов. В частности, построена обширная теория т. н. $A$-с к а л я р н ы х о п ер а т о р о в, применимая ко многим классам операторов, не укладывающихся в классическую спектральную теорию. Для успешного использования Ф.и. имеют значения т. н. теоремы об отображении спектра:



\[

\sigma(f(T))=f(\sigma(T)), f \in A

\]



для всех перечисленных выше Ф. и. такие теоремы доказаны (после надлежащего осмысления правой части формулы).



Если алгебра $A$ содержит мелкие разбиения единицы (вапр., $A=C^{\infty}$ ), то $A$-Ф. и. позволяет построить локальный спектральный анализ и, в частности, доказать существование нетривиальных инвариантных подпространств оператора $T$ (если $\sigma(T)$ - не одноточечное множество); пример - оператор $T$ (в банаховом пространстве), спектр которого лежит на гладкой кривой $\gamma$, и $S_{0} \ln +\ln +\delta(r) d r<\infty$, где $\delta(r)=$ $=\max \{\|R(\lambda ; T)\|: \operatorname{dist}(\lambda, \gamma) \geqslant r Y\}$. Следствием локального анализа является и теорема Шилова об идемпотентах [2].



Лum.. [1] Д а н ф о р д Н., ІІ в а р ц Д ж., Јинейные операторы, пер. с англ., ч. 1, М., 1962, ч. 3, М., 1974, [2] Б у р б ак и Н., Спектральная теория, пер. с франц., М., 1972, [3] W ae 1 b r o e c k L., Etude spectrale des algè bres complètes, Brux., 1960; [4] T a y l o r J. L., "Acta math.», 1970, v. 125, № 1-2, p. 1 - 38; [5] Д ы н ь кинн'Е. М., «Записки науч. семин. ЈОМИ», 1972, T. 30, c. 33-39; [6] v o n N e u m a n n J., «Math. Nachr.», $1950 / 51$, Bd 4, S. 258-81, [7] С е ке ф а ль в и-н а д ь Б. Ф о я ш Ч., Гармонический анализ операторов в гильбертовом пространстве, пер. с франц., М., 1970, [8] П Е л л е р В. В.,



\begin{center}

\includegraphics[max width=\textwidth]{2023_07_09_3ad187aa3f1a83972325g-1}

\end{center}



räl., F o i a $\$$ C., Theory of generalized spectral operators «. Ү.- L., 1968; [10] Л ю б и ч Ю. И., М а ца е в В. И., с к и й Я., Операторное исчисление, пер. с польск., М., 1956; [12] М а с л о в В. П., Операторные методы, М., 1973.



\includegraphics[max width=\textwidth, center]{2023_07_09_3ad187aa3f1a83972325g-1(1)}

ношение $R$ на множестве $A$, удовлетворяющее соотношению $R^{-1} \circ R \sqsubseteq \Delta$, где $\Delta$ - диагональ множества $A$. Әто означает, что из $(a, b) \in R$ и $(a, c) \in R$ следует, что $b=c$, т. е. каждому $a \in A$, стоящему на первом месте, соответствует не более одного $b \in A$ на втором месте в паре из $R$. Таким образом, отношение $R$ определяет (быть может, не всюду определенную) функцию на множестве $A$. При выполнении соотношения $R^{-1} \circ R=\Delta$ эта функция всюду определена и взаимно однозначна. О. А. Иванова.



ФУНКЦИОНАЛЬНОЕ УРАВНЕНИЕ - уравнение (линейное или нелинейное), в к-ром неизвестным является элемент какого-либо банахова пространства, конкретного (функционального) или абстрактного, т. е. уравнение вида



\[

P(x)=y,

\]



где $\boldsymbol{P}(x)$ - нек-рый, вообще говоря, нелинейный оператор, переводящий элементы пространства $X$ типа $B$ в элементы пространства $Y$ того же типа. Если Ф. у. содержит еще и числовой (или общий функциональный) параметр $\lambda$, то вместо (1) пишут



\[

P(x ; \lambda)=y,

\]



где $x \in X, y \in Y, \lambda \in \Lambda, \Lambda$ - пространство параметров. Уравнениями вида (1) являются конкретные или абстрактные диффференциальные уравнения, обыкновенные и с частными производными, интегральные уравнения, интегро-диффференциальные уравнения, функционально-дифференциальные и более сложные уравнения математич. анализа, а также системы алгебраич. уравнений, конечные и бесконечные, уравнения в конечных разностях и др.



В линейном случае рассматриваются Ф. у. 1-го рода $A x=y$ и 2 -го рода $x-\lambda A x=y$, где $A$ - линейный опө ратор из $X$ в $Y$, а $\lambda$ - параметр. При әтом формально Ф. у. 2-го рода может быть записано в виде уравнения 1-го рода $\operatorname{Tx}=y(T=I-\lambda A)$. Однако выделение тождественного оператора $I$ оказывается целесообразным, так как оператор $A$ может обладать лучшими свойствами, чем оператор $T$, что позволяет полнее исследовать. рассматриваемое уравнение.



Ф. у. рассматриваются также и в др. пространствах, напр. в пространствах, нормированных әлементами полуупорядоченных пространств.



Если решения Ф. у. являются элементами пространства операторов, то такие Ф. у. наз. о п е р а т о рны м и у р а в не и я м и (о. у.), ковкретными или абстрактными. Здесь также могут быть алгебраич. о. у., линейные и нелинейные, дифференциальные, интегральные и другие о. у. Напр., пусть в нормированном кольце $[X]=[X \rightarrow X]$ линейных операторов, переводящих пространство $X$ типа $B$ в себя, рассматривается обыкновенное дифференциальное уравнение на бесконечном промежутке $0 \leqslant \lambda \leqslant \infty$ :



\[

\frac{d x}{d \lambda}=-A x(=-x A), x(0)=x_{0},

\]



где $A, x_{0} \in[X], x(\lambda)$ - абстрактная функция со значениями в банаховом пространстве $[X]$. Это уравнение является простейшим абстрактным линейным дифференциальным о. у., оно получается, напр., из применения прямого метода вариации параметра к построению операторов вида $P(A)=e-A \lambda, 0 \leqslant \lambda \leqslant \infty$, в частности проекторов $P(A)\left([P(A)]^{2}=P(A)\right)$ с единичной нормой. Проекторы вида $P(A), P(A C)$ и $P(C A), C \in[X]$, применяются, напр., при построении прямым методом вариации параметра явных и неявных псевдообратных операторов и псевдорешений линейных Ф. У., а также собственных элементов (собственных подпространств) оператора $A$. Сведение различных задач к задачæм для уравнения (2) и др. является весьма удобным при разработке приближенных методов решения. Интерес представляют также о. у. вида



\[

\frac{d x}{d \lambda}=A(\lambda) x+F(\lambda)(=x A(\lambda)+F(\lambda)), x(0)=x_{0},

\]





\end{document}